\section{讨论}
\label{sec:discussion}

\subsection{三组试验形成的证据链}

三组试验在不同层级检验模型。算例组 1 检验计算的气液场能否复现单次事件的完整序列。
算例组 2 检验这些场量能否在竖管直径与气囊体积改变时选对分支。算例组 3 则检验交汇室
和下游边界参与瞬变后，降阶表述是否仍有解释力。综合结果表明，守恒并输运气相足以恢复
三套装置中的主要机理，但选对分支并不等同于事件时刻、界面速度和压力峰均准确。

算例组 1 中，宽塔与窄塔分属不同试验分支，模型正确区分了二者。宽塔计算最大水位为
$0.66L$，实测为 $0.63L$；界面速度为 $0.41\sqrt{gD}$，实测为
$0.39\sqrt{gD}$。压力平台也接近实测。这些结果说明，解析气囊体积和静水阻力正确控制了
整体响应。残余误差来自拓扑：试验中存在带排液壁膜的连贯气囊，而当前网格尺度闭合把
井内气体作为弥散相输运，因此没有复现水面缓升，并把压力衰减推迟至界面突破之后。

窄塔暴露出另一类误差。模型复现了喷发前驻位（$0.84L$ 对 $0.82L$）、压力平台、到达
井口及水面速度（$0.39\sqrt{gD}$ 对 $0.44\sqrt{gD}$）。然而，界面速度仅为
$0.71\sqrt{gD}$，实测为 $1.43\sqrt{gD}$。水面到达井口偏早（$t^{*}=3.70$ 对
$3.90$），界面突破却偏晚（$t^{*}=4.44$ 对约 $3.85$）。因此，误差是气体释放序列被
拉长，而不是全部事件统一平移。这一特征把问题定位到井口供气和气囊供压气核，而不是仅仅
定位到初始气囊到达。

算例组 2 说明该误差如何影响分支判别。模型在 7 个 Series-B 录像工况中判对 5 个，在
63 个完整参数扫描构型中判对 39 个。它复现了喷发端元、全部宽竖管和小气囊不喷发端元，
以及管端压力的实测分支差异。24 个扫描分歧均为试验喷发而计算水位未及井口，没有计算假
喷发。该单向模式与过渡带附近净水柱抬升不足一致。在风险筛查语境中不宜称其为``保守''，
因为漏报喷发对安全而言并不保守。

算例组 3 提供了互补的边界条件检验。仅改变下游条件即可使 A2 与 B3 的结局翻转，模型
复现了这一翻转，也使 B3 喷高落在试验压力--高度趋势上。冻结波速低估了短时压力峰，
计算为 $30.84$~kPa，实测为 $55.03$~kPa。C9 的计算首峰为 $11.49$~kPa、时刻
$0.54$~s，实测为 $10.69$~kPa、$0.50$~s。计算振荡周期偏短，后续气囊驱动喷发则超出
复合一维表述范围。因此，降阶交汇室保留了首次瞬变响应，但未保留全部后期气囊动力学。

\subsection{降阶模型能够可靠解析的量}

最重要的结果是：模型无需输入气体释放时程即可恢复分支物理。各组计算中，气囊压力均由
守恒气体质量与计算体积演化，进入竖井的气体输运和水柱位移随耦合解产生。模型因而能够
响应竖管直径、气囊体积和下游边界变化，并区分平缓排气与喷发。液相模型加规定压力或流量
可以拟合选定迹线，却不能提供同样独立的分支检验。

对比还区分了在一维分辨率下较稳健和较敏感的观测量。缓变压力平台、准静态水柱驻位、
最大水位及首次质量振荡峰，比界面突破、环状液膜运动和短时声学峰更可靠。前一类量主要
受气体存量、水柱惯性与静水平衡控制；后一类量则依赖结点和井壁处的局部三维交换与未解析
时间尺度。

\subsection{试验揭示的闭合局限}

首要局限是井口闭合。窄塔受驱区会启用管顶压力基准、气囊供压气核和 Wallis 型淹没限。
这些要素使用几何条件与文献尺度常数，而非逐工况标定；但它们仍是按流态触发的闭合，而非
统一的截面结点解。宽塔采用管底基准，窄塔则需要物理上的管顶连接。后续结点模型应传递
井口截面的竖向压力分布，而不是压缩为单一标量基准。

第二项局限是气穴传播与排气供给相互绑定。冻结的耗散气穴速度使初始到达落入试验散布，
但同一降低后的供给也减慢气囊供压气核并推迟突破。若把前沿速度与压差驱动的井口气体通量
分开闭合，两项约束才可能独立满足。

第三项局限是 T 型结点的体积中性逆向交换。该交换可在平衡排气时维持气囊压力，也是产生
B-H1 喷发所必需的机制；但在水库供水的宽竖管中，它抵消了进入气体本应造成的部分水柱
位移。直接删除交换虽能恢复宽竖管抬升，却同时破坏 B-H1 的压力重建。两分支因此需要以
当地逆流容量为基础的限流交换律，而不是强制气液体积相等。

另有两项局限来自维数降阶。弥散两流体状态不能维持宽塔中连贯泰勒气囊及其排液膜拓扑。
降阶交汇室和软化声学波速也不能解析 B3 最短压力峰及 C9 后期气囊到达和排放。这些误差
应通过有针对性的闭合或局部多维耦合解决，而不应通过逐试验调整全局参数掩盖。

\subsection{应用含义与范围}

现有证据支持将该框架用于机理研究，以及远离分支边界的系统尺度筛查。它能在一个计算
描述中跟踪气体存量、气囊压缩、分支选择以及整体压力与水柱响应。但现有证据不支持把它
直接作为边界附近的独立安全判别器。算例组 2 的误差是漏报喷发，B3 局部峰压还对声学
表述敏感。

本文精度只适用于冻结的数值配置和已说明的流态条件结点闭合，并不证明模型可直接推广至
任意井形、井盖、管网相互作用或原型雷诺数与韦伯数。下一步应在不修改闭合参数的前提下，
对独立装置进行验证。优先量测量应包括井口同步气体通量、液体回流、气囊压力和界面拓扑。
这些数据可区分交换律误差与竖井卷吸误差，使下一轮闭合改进具有可证伪性。

\section{结论}
\label{sec:conclusions}

本文用三组已发表喷发试验评估了守恒气相的一维混合流框架。竖井与水平管道采用耦合两流
体状态，气体释放由计算产生而非预设。通风竖井试验中，模型正确选择不喷发与喷发分支，
并复现主要压力和水柱尺度。窄塔对比显示释放序列被拉长：气体到达偏早，界面突破偏晚，
计算界面速度约为实测的一半。

水平管--竖管试验中，模型在 7 个录像工况中判对 5 个，在 63 个参数扫描构型中判对
39 个。24 个扫描误差全部是过渡带附近的实测喷发漏报，没有假喷发。交汇室试验中，模型
复现 A2/B3 分支翻转及 C9 首次压力瞬变，但低估了 B3 的短时压力峰，且未表示后续气囊
驱动事件。

综合证据把系统尺度能力与当前局限清楚分开。守恒并输运气相可以恢复主要分支机理、整体
压力尺度与水柱响应。剩余误差集中在井口交换、连贯液膜拓扑以及声学与交汇室降阶。因此，
该框架适合机理评估和远离边界的有条件筛查。若要用于临界喷发的安全判别，还必须建立统一
且受通量限制的结点与卷吸闭合。
